% GENERATED FILE. Edit canonical lessons, then run npm run generate:books.
% canonical-source-hash: fnv1a64:40595e903c1f938d
% canonical-lessons: GE-C63-nachbar, GE-C63-gast, GE-C63-enkel, GE-C63-schueler, GE-C63-schuelerin

\chapter{Neighbour, Guest, Grandson, Pupil, Pupil}
\label{ch:ge-neighbour-guest-grandson-pupil-pupil}
\hlchaptermodality{63}

\begin{chapteropening}
\textbf{By the end of this chapter:} \emph{I can say a neighbour, a guest, a grandson, a pupil and a pupil.}

Complete the last lesson of chapter 63: I can say a neighbour, a guest, a grandson, a pupil and a pupil.
\end{chapteropening}

\section[der Nachbar]{der Nachbar — a neighbour}

\label{lesson:GE-C63-nachbar}



\begin{quote}
Before the new one: say the German for an uncle, then the German for an aunt.
\end{quote}

\subsection*{You'll want to know: der Nachbar}
\textbf{der Nachbar} — \textquotedblleft{}a neighbour\textquotedblright{}.

\begin{cousinweb}
\textbf{Nachbar} $\leftarrow$ Middle High German \emph{nāchbūre} (EtymDB 2.1), \textquotedblleft{}the one who dwells near\textquotedblright{}: the word for \textquotedblleft{}near\textquotedblright{} + an old word for \textquotedblleft{}dweller\textquotedblright{}. English \emph{neighbour} is built the same way.
\end{cousinweb}

\begin{grammarlens}[title={What you've built}]
One more word for this chapter.
\end{grammarlens}

\subsection*{Your turn}
\begin{itemize}
\raggedright
  \item \emph{Say it:} \emph{der Nachbar}
  \item \emph{Say it:} \emph{der Nachbar}, once more
  \item \emph{Say it:} say \emph{der Nachbar}
  \item \emph{Recall:} say the German for an uncle, then the German for an aunt, then say \emph{der Nachbar} again
\end{itemize}

\begin{tcolorbox}[breakable,colback=teal!4,colframe=teal!35!black,title={Before you move on}]
What does der Nachbar mean? (\textquotedblleft{}A neighbour\textquotedblright{}.) Say it once more.
\end{tcolorbox}
\section[der Gast]{der Gast — a guest}

\label{lesson:GE-C63-gast}



\begin{quote}
Before the new one: say the German for an aunt, then the German for a neighbour.
\end{quote}

\subsection*{You'll want to know: der Gast}
\textbf{der Gast} — \textquotedblleft{}a guest\textquotedblright{}.

\begin{cousinweb}
\textbf{Gast} $\leftarrow$ Old High German \emph{gast} $\leftarrow$ Proto-Germanic \emph{*gastiz}, \textquotedblleft{}guest, stranger\textquotedblright{} — inherited, the same word as English \emph{guest}. (EtymDB 2.1)
\end{cousinweb}

\begin{grammarlens}[title={What you've built}]
One more word for this chapter.
\end{grammarlens}

\subsection*{Your turn}
\begin{itemize}
\raggedright
  \item \emph{Say it:} \emph{der Gast}
  \item \emph{Say it:} \emph{der Gast}, once more
  \item \emph{Say it:} say \emph{der Gast}
  \item \emph{Recall:} say the German for an aunt, then the German for a neighbour, then say \emph{der Gast} again
\end{itemize}

\begin{tcolorbox}[breakable,colback=teal!4,colframe=teal!35!black,title={Before you move on}]
What does der Gast mean? (\textquotedblleft{}A guest\textquotedblright{}.) Say it once more.
\end{tcolorbox}
\section[der Enkel]{der Enkel — a grandson}

\label{lesson:GE-C63-enkel}



\begin{quote}
Before the new one: say the German for a neighbour, then the German for a guest.
\end{quote}

\subsection*{You'll want to know: der Enkel}
\textbf{der Enkel} — \textquotedblleft{}a grandson\textquotedblright{}.

\begin{cousinweb}
\textbf{Enkel} $\leftarrow$ Middle High German \emph{enikel} $\leftarrow$ Old High German \emph{eniklīn}, a diminutive of \emph{ano}, \textquotedblleft{}grandfather\textquotedblright{} (EtymDB 2.1) — a grandson was once a \textquotedblleft{}little grandfather\textquotedblright{}.
\end{cousinweb}

\begin{grammarlens}[title={What you've built}]
One more word for this chapter.
\end{grammarlens}

\subsection*{Your turn}
\begin{itemize}
\raggedright
  \item \emph{Say it:} \emph{der Enkel}
  \item \emph{Say it:} \emph{der Enkel}, once more
  \item \emph{Say it:} say \emph{der Enkel}
  \item \emph{Recall:} say the German for a neighbour, then the German for a guest, then say \emph{der Enkel} again
\end{itemize}

\begin{tcolorbox}[breakable,colback=teal!4,colframe=teal!35!black,title={Before you move on}]
What does der Enkel mean? (\textquotedblleft{}A grandson\textquotedblright{}.) Say it once more.
\end{tcolorbox}
\section[der Schüler]{der Schüler — a pupil (a boy)}

\label{lesson:GE-C63-schueler}



\begin{quote}
Before the new one: say the German for a guest, then the German for a grandson.
\end{quote}

\subsection*{You'll want to know: der Schüler}
\textbf{der Schüler} — \textquotedblleft{}a pupil (a boy)\textquotedblright{}.

\begin{cousinweb}
\textbf{Schüler} $\leftarrow$ Middle High German \emph{schuolære}, built on the word for \textquotedblleft{}school\textquotedblright{}, Old High German \emph{scuola} $\leftarrow$ Latin \emph{schola} (EtymDB 2.1) — a borrowed stem with the German ending \emph{-er}.
\end{cousinweb}

\begin{grammarlens}[title={What you've built}]
One more word for this chapter.
\end{grammarlens}

\subsection*{Your turn}
\begin{itemize}
\raggedright
  \item \emph{Say it:} \emph{der Schüler}
  \item \emph{Say it:} \emph{der Schüler}, once more
  \item \emph{Say it:} say \emph{der Schüler}
  \item \emph{Recall:} say the German for a guest, then the German for a grandson, then say \emph{der Schüler} again
\end{itemize}

\begin{tcolorbox}[breakable,colback=teal!4,colframe=teal!35!black,title={Before you move on}]
What does der Schüler mean? (\textquotedblleft{}A pupil (a boy)\textquotedblright{}.) Say it once more.
\end{tcolorbox}
\section[die Schülerin]{die Schülerin — a pupil (a girl)}

\label{lesson:GE-C63-schuelerin}



\begin{quote}
Before the new one: say the German for a grandson, then the German for a pupil.
\end{quote}

\subsection*{You'll want to know: die Schülerin}
\textbf{die Schülerin} — \textquotedblleft{}a pupil (a girl)\textquotedblright{}.

\begin{cousinweb}
\textbf{Schülerin} is \emph{Schüler} + \emph{-in}, the ending that names a woman or girl (EtymDB 2.1). The stem goes back to Latin \emph{schola}, \textquotedblleft{}school\textquotedblright{}.
\end{cousinweb}

\begin{grammarlens}[title={What you've built}]
One more word for this chapter.
\end{grammarlens}

\subsection*{Your turn}
\begin{itemize}
\raggedright
  \item \emph{Say it:} \emph{die Schülerin}
  \item \emph{Say it:} \emph{die Schülerin}, once more
  \item \emph{Say it:} say \emph{die Schülerin}
  \item \emph{Recall:} say the German for a grandson, then the German for a pupil, then say \emph{die Schülerin} again
\end{itemize}

\begin{tcolorbox}[breakable,colback=teal!4,colframe=teal!35!black,title={Before you move on}]
What does die Schülerin mean? (\textquotedblleft{}A pupil (a girl)\textquotedblright{}.) Say it once more.
\end{tcolorbox}
